\clearpage
\renewcommand{\sectioncolor}{bio}
\section{What to actually do}\label{sec:do}
\markboth{What to do}{}

\subsection{Where the original worry is legitimate}

I do not want to talk you entirely out of it. A residue is real, and it is worth naming precisely so
that it can be worked around rather than worried about.

\begin{center}\small
\begin{tabular}{@{}>{\raggedright\arraybackslash}p{5.3cm}>{\raggedright\arraybackslash}p{11.4cm}@{}}
\toprule
\rowcolor{crimson!11}
\textbf{\color{crimson}Residual limit} & \textbf{\color{crimson}What it actually costs}\\
\midrule
A speed and reliability tax &
 Unintuitive mathematics takes longer to invent, is harder to check, and is taught badly. The
 parallel postulate cost roughly two millennia, and 39 more years between the symbolic escape and
 the drawable model. \textbf{This is an embodiment effect and it is expensive.}\\
Ceilings unrelated to embodiment &
 Gödel-type incompleteness, and computational intractability, bind any mind. Many-body quantum
 simulation is exponentially hard for reasons that have nothing to do with how brains work.\\
\rowcolor{crimson!6}
A selection effect &
 We may be systematically failing to notice mathematics we never invented. As stated this is
 unfalsifiable; hold it lightly, but it is not nothing.\\
\bottomrule
\end{tabular}
\end{center}

\subsection{Four things one person can do}

\begin{movebox}[title={1. Do not try to invent ``the mathematics of the cell''}]
Historically, mathematics that works for a new domain is invented by people forced to compute
\emph{one specific blocked thing}.
\begin{center}
\begin{tikzpicture}[font=\scriptsize,
  h/.style={rounded corners=3pt,draw=bio,line width=0.85pt,fill=bio!8,inner sep=6pt,
            text width=3.15cm,align=center,minimum height=1.3cm}]
 \node[h] at (0,0)     {\textbf{Fourier}\\{\scriptsize had a heat equation}};
 \node[h] at (3.9,0)   {\textbf{Heaviside}\\{\scriptsize had a telegraph cable}};
 \node[h] at (7.8,0)   {\textbf{Feynman}\\{\scriptsize had a QED calculation}};
 \node[h] at (11.7,0)  {\textbf{Feinberg}\\{\scriptsize had a reaction network}};
 \node[h,draw=crimson,fill=crimson!8] at (15.6,0) {\textbf{you}\\{\scriptsize need one blocked calculation}};
\end{tikzpicture}
\end{center}
\vspace{-2pt}
\textbf{Pick one calculation you cannot currently do, and invent only what it needs.} Generality is
what you discover afterwards, by stripping the biology out (the pasted text's step 7 --- which is
correct, but comes seventh for a reason).
\end{movebox}

\begin{movebox}[title={2. Run the metaphor audit as a daily habit}]
For any formalism you meet, run Part VI's four questions from p.~436 --- what blend makes this true,
what relations does it presuppose, what do the symbols mean, where is it grounded --- and then add a
fifth that belongs to this dossier:
\begin{center}
\begin{tikzpicture}
\node[rounded corners=4pt,fill=bio!12,draw=bio,line width=1.1pt,inner sep=10pt,
      text width=15.2cm,align=center,font=\itshape]
 {\textbf{Which grounding in this formalism is the one that will break at $10^{-9}$\,m?}};
\end{tikzpicture}
\end{center}
\end{movebox}

\begin{movebox}[title={3. Learn the three languages built for post-intuitive domains}]
\begin{center}\small
\begin{tabular}{@{}>{\raggedright\arraybackslash}p{4.0cm}>{\raggedright\arraybackslash}p{5.6cm}>{\raggedright\arraybackslash}p{6.5cm}@{}}
\toprule
\rowcolor{bio!13}
\textbf{\color{bio}Language} & \textbf{\color{bio}What it replaces} & \textbf{\color{bio}Why it matters here}\\
\midrule
\textbf{Category theory} & composition, not collection &
 cells are relational and compositional, not set-like\\
\textbf{Sheaves} & local data, global obstruction &
 context-dependent identity --- \S\ref{ss:context}\\
\textbf{A diagrammatic calculus}\newline{\scriptsize string diagrams, ZX} & re-grounds computation in hand and eye &
 Move 3, the highest-yield move in the history of this problem\\
\bottomrule
\end{tabular}
\end{center}
\end{movebox}

\begin{movebox}[title={4. Treat instruments as mathematics-generators, not just data sources}]
Every new sensorimotor loop into a new scale has produced new mathematics within roughly a decade.
If you are choosing between building a theory for data that exists and building a way to see
something nobody has seen, the historical record favours the second more often than
mathematicians like to admit.
\end{movebox}

\subsection{The one-sentence answer}

\begin{center}
\begin{tikzpicture}
\node[rounded corners=5pt,fill=bio!92!black,text=white,inner sep=13pt,text width=15.8cm,align=center]
 {\large You cannot escape embodiment, and you do not need to.\\[7pt]
  \normalsize\normalfont
  \textbf{Extend the body} with instruments; \textbf{re-ground the abstraction} with diagrammatic
  calculi; \textbf{externalise the bookkeeping} to machines.\\[5pt]
  That three-part move has worked every time it has been tried.};
\end{tikzpicture}
\end{center}

\subsection{The question worth giving your students}

The pasted text ends by proposing a question for two of your mentees. Their names have been removed
here, but the question itself is the best line in the piece and deserves sharpening:

\begin{center}
\begin{tikzpicture}[font=\footnotesize,
  q/.style={rounded corners=4pt,draw=#1,line width=1.2pt,fill=#1!7,inner sep=10pt,
            text width=7.45cm,align=left,minimum height=3.1cm}]
 \node[q=crimson] (a) at (-4.35,0)
   {\textbf{\color{crimson}Not}\\[5pt]
    \textit{\large ``How much mathematics can you memorise?''}\\[5pt]
    {\scriptsize Tests retrieval. Predicts nothing about research capability.}};
 \node[q=bio] (b) at (4.35,0)
   {\textbf{\color{bio}But}\\[5pt]
    \textit{\large ``When the mathematics you know fails to explain nature, can you find the
    assumption that is doing the damage?''}\\[5pt]
    {\scriptsize Tests the metaphor audit --- Move 1, and the one skill the whole method rests on.}};
\end{tikzpicture}
\end{center}
\vspace{2pt}
\noindent The pasted version asks whether a student can \emph{``invent a better representation.''}
That is Move 3, and it is the right ambition --- but it is the second step. \textbf{You cannot invent
a better representation until you have located which embodied assumption the current one is smuggling
in.} Newton needed calculus because he had first noticed that the available mathematics could not
express instantaneous change; Connes needed non-commutative geometry because he had first noticed
that ``space is a set of points'' was the offending premise.

Teach the diagnosis before the cure.
